% META: {"id": "TCOMPL-ECHANTILLONNAGE-QCM", "chapitre": "TCOMPL-ECHANTILLONNAGE", "type_objet": "qcm", "genere_depuis": "chapitres/TCOMPL-ECHANTILLONNAGE/qcm/TCOMPL-ECHANTILLONNAGE-QCM.json", "status": "generated"}
% Fichier genere par scripts/build_qcm_tex.py — ne pas editer a la main.

\section*{\textcolor{chapcolor}{\MakeUppercase{Echantillonnage et estimation}}}

\begin{center}
\textit{Pour chaque question, une seule réponse est exacte.}
\end{center}

\begin{enumerate}

\item[] \textbf{Capacité C1}

\item \textbf{[Q1]} Le nombre $X$ de succès dans $n$ repetitions indépendantes d'une même epreuve de probabilité de succès $p$ suit :
  \begin{enumerate}[label=\Alph*.]
    \item la loi binomiale $\mathcal B(n,p)$
    \item une loi uniforme
    \item une loi exponentielle
    \item une loi normale dans tous les cas
  \end{enumerate}

\item[] \textbf{Capacité C2}

\item \textbf{[Q2]} Quelle relation du triangle de Pascal permet de calculer $\binom{n}{k}$ ?
  \begin{enumerate}[label=\Alph*.]
    \item $\binom{n}{k}=\binom{n-1}{k-1}+\binom{n-1}{k}$
    \item $\binom{n}{k}=n^k$
    \item $\binom{n}{k}=\binom{n-1}{k-1}-\binom{n-1}{k}$
    \item $\binom{n}{k}=\binom{k}{n}$
  \end{enumerate}

\item[] \textbf{Capacité C3}

\item \textbf{[Q3]} Si $X$ suit $\mathcal B(3,1/2)$, alors $P(X=2)$ vaut :
  \begin{enumerate}[label=\Alph*.]
    \item $1/8$
    \item $3/8$
    \item $1/2$
    \item $3/4$
  \end{enumerate}

\item[] \textbf{Capacité C4}

\item \textbf{[Q4]} Si $X$ suit $\mathcal B(4,1/2)$, alors $P(1\leq X\leq3)$ vaut :
  \begin{enumerate}[label=\Alph*.]
    \item $1/4$
    \item $1/2$
    \item $7/8$
    \item $3/4$
  \end{enumerate}

\item[] \textbf{Capacité C5}

\item \textbf{[Q5]} Quelle expression Python renvoie la moyenne d'un échantillon de $n$ variables de Bernoulli de paramètre $p$ ?
  \begin{enumerate}[label=\Alph*.]
    \item sum(range(n)) / p
    \item sum(random() < p for \_ in range(n))
    \item random() < p / n
    \item sum(random() < p for \_ in range(n)) / n
  \end{enumerate}

\item[] \textbf{Capacité C6}

\item \textbf{[Q6]} Soit $X$ suivant $\mathcal B(3;\tfrac12)$. Le calcul direct $E(X)=\sum_{k=0}^{3} k\,P(X=k)$ donne :
  \begin{enumerate}[label=\Alph*.]
    \item $\tfrac12$
    \item $\tfrac32$
    \item $3$
    \item $\tfrac34$
  \end{enumerate}

\item[] \textbf{Capacité C7}

\item \textbf{[Q7]} Si $X$ suit la loi uniforme sur $\{1;2;\dots;n\}$, son espérance vaut :
  \begin{enumerate}[label=\Alph*.]
    \item $\tfrac1n$
    \item $\tfrac n2$
    \item $\tfrac{n+1}2$
    \item $n$
  \end{enumerate}

\end{enumerate}

% NEXUS-QCM-TEACHER-ONLY-BEGIN
\ifnxVersionProfesseur
\clearpage
\section*{Clé de correction — réservée au professeur}

\begin{center}
\begin{tabular}{lll}
\hline
Question & Capacité & Réponse exacte \\
\hline
Q1 & C1 & \textbf{A} \\
Q2 & C2 & \textbf{A} \\
Q3 & C3 & \textbf{B} \\
Q4 & C4 & \textbf{C} \\
Q5 & C5 & \textbf{D} \\
Q6 & C6 & \textbf{B} \\
Q7 & C7 & \textbf{C} \\
\hline
\end{tabular}
\end{center}

\subsection*{Diagnostic des réponses erronées}

\begin{itemize}
  \item \textbf{Q1}
  \begin{itemize}
    \item \textbf{B} — Une loi uniforme attribuerait la même probabilité à toutes les valeurs, ce que le nombre de succès ne fait pas. Le nombre de succès dans n répétitions indépendantes de même probabilité p suit la loi binomiale B(n,p). \emph{Renvoi : C1.}
    \item \textbf{C} — La loi exponentielle modélise un temps d'attente continu, pas un nombre discret de succès. Le nombre de succès dans n répétitions indépendantes de même probabilité p suit la loi binomiale B(n,p). \emph{Renvoi : C1.}
    \item \textbf{D} — Une approximation normale peut parfois être utilisée, mais elle n'est ni exacte ni valable dans tous les cas. Le nombre de succès dans n répétitions indépendantes de même probabilité p suit la loi binomiale B(n,p). \emph{Renvoi : C1.}
  \end{itemize}
  \item \textbf{Q2}
  \begin{itemize}
    \item \textbf{B} — La puissance $n^k$ compte des listes ordonnées avec répétition, pas des sous-ensembles de taille $k$. La relation de Pascal est C(n,k)=C(n-1,k-1)+C(n-1,k), obtenue selon que l'élément distingué est choisi ou non. \emph{Renvoi : C2.}
    \item \textbf{C} — Les deux classes disjointes se réunissent par addition et non par soustraction. La relation de Pascal est C(n,k)=C(n-1,k-1)+C(n-1,k), obtenue selon que l'élément distingué est choisi ou non. \emph{Renvoi : C2.}
    \item \textbf{D} — Inverser $n$ et $k$ rend généralement le coefficient nul ou non défini et ne traduit pas la relation de Pascal. La relation de Pascal est C(n,k)=C(n-1,k-1)+C(n-1,k), obtenue selon que l'élément distingué est choisi ou non. \emph{Renvoi : C2.}
  \end{itemize}
  \item \textbf{Q3}
  \begin{itemize}
    \item \textbf{A} — La valeur $1/8$ est la probabilité d'une configuration précise de deux succès et oublie les trois positions possibles. Pour X suivant B(3,1/2), P(X=2)=C(3,2)(1/2)\textasciicircum{}2(1/2)=3/8. \emph{Renvoi : C3.}
    \item \textbf{C} — La valeur $1/2$ ne résulte pas de la formule binomiale et confond le paramètre $p$ avec la probabilité de deux succès. Pour X suivant B(3,1/2), P(X=2)=C(3,2)(1/2)\textasciicircum{}2(1/2)=3/8. \emph{Renvoi : C3.}
    \item \textbf{D} — La valeur $3/4$ oublie le facteur $(1/2)^3$ associé aux trois résultats de chaque configuration. Pour X suivant B(3,1/2), P(X=2)=C(3,2)(1/2)\textasciicircum{}2(1/2)=3/8. \emph{Renvoi : C3.}
  \end{itemize}
  \item \textbf{Q4}
  \begin{itemize}
    \item \textbf{A} — La valeur $1/4$ ne somme pas les probabilités des trois valeurs 1, 2 et 3 de l'intervalle. Pour X suivant B(4,1/2), P(1<=X<=3)=(4+6+4)/16=14/16=7/8. \emph{Renvoi : C4.}
    \item \textbf{B} — La valeur $1/2$ ne tient compte que du coefficient central et omet les cas $X=1$ et $X=3$. Pour X suivant B(4,1/2), P(1<=X<=3)=(4+6+4)/16=14/16=7/8. \emph{Renvoi : C4.}
    \item \textbf{D} — La valeur $3/4$ retranche à tort quatre issues sur seize, alors que seules $X=0$ et $X=4$ sont exclues. Pour X suivant B(4,1/2), P(1<=X<=3)=(4+6+4)/16=14/16=7/8. \emph{Renvoi : C4.}
  \end{itemize}
  \item \textbf{Q5}
  \begin{itemize}
    \item \textbf{A} — La somme des indices ne simule aucune variable de Bernoulli et la division par $p$ n'a pas de sens statistique. La moyenne d'un échantillon de n Bernoulli simulées est la somme des indicatrices divisée par n. \emph{Renvoi : C5.}
    \item \textbf{B} — Cette expression renvoie le nombre de succès, mais oublie la division par la taille $n$ nécessaire pour obtenir la moyenne. La moyenne d'un échantillon de n Bernoulli simulées est la somme des indicatrices divisée par n. \emph{Renvoi : C5.}
    \item \textbf{C} — Cette expression ne simule qu'une seule épreuve et modifie en plus sa probabilité en utilisant $p/n$. La moyenne d'un échantillon de n Bernoulli simulées est la somme des indicatrices divisée par n. \emph{Renvoi : C5.}
  \end{itemize}
  \item \textbf{Q6}
  \begin{itemize}
    \item \textbf{A} — La valeur $1/2$ est la probabilité de succès $p$, pas une moyenne du nombre de succès. Le calcul direct donne $E(X)=0\times\tfrac18+1\times\tfrac38+2\times\tfrac38+3\times\tfrac18=\tfrac{12}8=\tfrac32$, soit $np$. \emph{Renvoi : C6.}
    \item \textbf{C} — La valeur $3$ est le nombre de répétitions $n$ : ce serait l'espérance si chaque répétition était un succès certain. Le calcul direct donne $E(X)=\tfrac{12}8=\tfrac32$, soit $np$. \emph{Renvoi : C6.}
    \item \textbf{D} — La valeur $3/4$ est la variance $np(1-p)$, qui mesure la dispersion et non la moyenne. Le calcul direct donne $E(X)=\tfrac{12}8=\tfrac32$, soit $np$. \emph{Renvoi : C6.}
  \end{itemize}
  \item \textbf{Q7}
  \begin{itemize}
    \item \textbf{A} — La valeur $1/n$ est la probabilité de chacune des $n$ valeurs, pas la moyenne de ces valeurs. L'espérance vaut $\tfrac1n\sum_{k=1}^{n}k=\tfrac1n\times\tfrac{n(n+1)}2=\tfrac{n+1}2$. \emph{Renvoi : C7.}
    \item \textbf{B} — La valeur $n/2$ serait la moyenne des entiers de $0$ à $n$ ; ici la plus petite valeur est $1$, pas $0$. L'espérance vaut $\tfrac1n\times\tfrac{n(n+1)}2=\tfrac{n+1}2$. \emph{Renvoi : C7.}
    \item \textbf{D} — La valeur $n$ est la plus grande valeur prise par $X$, pas sa moyenne : une moyenne ne peut pas égaler le maximum quand plusieurs valeurs distinctes sont possibles. L'espérance vaut $\tfrac{n+1}2$. \emph{Renvoi : C7.}
  \end{itemize}
\end{itemize}
\fi
% NEXUS-QCM-TEACHER-ONLY-END
